\section{Introduction}
\label{sec:intro}

Optimizing the use of resources in manufacturing supply chains is crucial for boosting economic growth, improving social well-being, and minimizing environmental impact.
Integrating transport, inventory, and production processes is a fundamental characteristic of supply chains \autocite{Stadtler2015SupplyPlanning} because joint decisions will improve supply chain performance\autocite{Diaz-Madronero2017ADecisions}.

Procurement refers to the process of acquiring goods and services in a supply chain \autocite{Chopra2015SupplyOperation}.
One of the main goals of procurement is to obtain \ac{RM}, parts or \ac{FG} for a factory or distribution centre from a supplier.
As mentioned in \autocite{Fleischmann2015TransportDistribution}, the responsibility for transport planning usually lies with the supplier.
However, on some \acp{SC}s the manufacturer makes transport decisions of its supplies in, for example, the automotive industry.
Furthermore, an essential asset in the automotive supply chain is not always considered in traditional \ac{MRP}: returnable containers or \acfp{RTI}, which are used as tertiary packaging.
\acp{RTI} enable the storage and transportation of \ac{RM} and \ac{FG} through the supply chain.
Usually, the manufacturing plant supplies empty \acp{RTI} to its supplier, while the supplier delivers full \acp{RTI} with the \ac{FG} or \ac{RM}.
Therefore, during the procurement process, transport and inventory decisions should be integrated, including \ac{RTI} management.
Here, we focus on optimising the operational decisions made during procurement transport planning.

According to \autocite{Diaz-Madronero2014AChain}, using algorithm-based methods for procurement transport planning is more efficient than relying on manual approaches based on planners' intuition.
Manual planning can be short-sighted and lead to less efficient outcomes.
This becomes even more significant when optimizing procurement transport tasks that aim to reduce the number of trucks and inventory levels, especially under 3D loading constraints where the precise placement of containers within trucks must also be determined.
Furthermore, the capacity constraints used in the scientific literature on procurement transport planning are based mainly on freight volume or the total number of loaded product units {\autocite{Hashemi-Amiri2023IntegratedApproach} and \autocite{Diaz-Madronero2014AChain}.
	}
The current research seeks to address this gap by introducing a new \ac{MILP} model.
This model uniquely combines procurement transportation with 3D full-truck loading, inventory management, and \ac{RTI} management.
Hereafter, we refer to this model as the \ac{RTISC}, and to its volumetric relaxation, in which truck capacity is a volume limit, as the \ac{VRTISC}.
Additionally, a \ac{BP} procedure is developed to solve both, and it is compared against the \ac{CF} solved directly by Gurobi, a commercial solver.

In this paper, we present four main contributions:
(i) We formulate a novel \ac{MILP} model, the \ac{RTISC}, integrating procurement transport, \ac{RTI} management and inventory decisions with stack-based full-truck loading constraints.
(ii) We use \ac{DW} decomposition to reformulate the problem over 3D loading configurations per time period. We define a \ac{BP} schema with column generation over that mixed integer pricing problem and a relax-and-fix construction as primal heuristic, giving the \ac{RTISCBP}.
(iii) We validate the accuracy of the \ac{BP} procedure against the \ac{CF} on both the 3D loading problem and its volumetric relaxation.
On the \ac{VRTISC}, where the \ac{VRTISCCF} reaches a proved optimum, the \ac{VRTISCBP} matches it to within $0.0001\%$ (Section~\ref{sec:res-vol}).
{\color{red} Comentar difencias entre BP y CF para el 3d}.
(iv) We show that the \ac{VRTISC}, the volumetric representation of truck capacity used in the procurement transport literature, yields plans that are cheaper on paper but cannot be loaded in practice, nor repaired by re-packing their own quantities.
{\color{red} Añadir algo vs company procedures}.

Readers are referred to the typology by \autocite{Wascher2007AnProblems} for a global view of container loading problems, where the different dimensions of the loading problem are described.
Furthermore, \autocite{Silva2019ExactProblems} summarise the exact methods used in the literature to solve container loading problems, and \autocite{Zhao2016INTERNATIONALAlgorithms} explore other solution approaches in the literature.
Readers are also referred to the literature reviews of \autocite{Mosca2019IntegratedReview} and \autocite{Diaz-Madronero2015ADecisions}, which focus more on integrating inventory decisions and transport.
Here, we jointly address container loading, procurement transport and inventory decisions.

The paper is organised as follows.
Section~\ref{sec:lit} provides a literature review related to the problem under study.
Section~\ref{sec:problem} presents the problem description.
Section~\ref{sec:model} describes the \ac{RTISC} \ac{MILP} formulation.
Sections~\ref{sec:decomp}--\ref{sec:branching} develop the \ac{BP} solution approach, the \ac{RTISCBP}, together with its pricing subproblem and bounding arguments.
Section~\ref{sec:results} reports the computational experiments.
Finally, Section~\ref{sec:conclusions} concludes and puts forward further research lines.
